% ==============================================================================
% SECTION 6: LIMITATIONS AND FUTURE WORK
% ==============================================================================

\section{Limitations and Future Work}

\subsection{Limitations}

The mandatory architectural bail-out protocol delivers strong safety, liveness, and
accountability guarantees under the structural invariants established in
Section~\ref{sec:guarantees}.  Nevertheless, four constraints bound the current
instantiation and must be addressed before deployment in adversarial or
high-throughput operational environments.

% ------------------------------------------------------------------------------
\subsubsection{Human Response Latency}

The protocol relies on a human principal as the final trust arbiter for any
escalated trigger.  Human cognitive response latency introduces a fundamental
bottleneck: mean response times for novel situations range from 200--500\,ms for
initial evaluation, with an additional 50--150\,ms for signal composition and
transmission~\cite{card1980psychology}.  In high-frequency trading engines,
autonomous vehicle control loops, or real-time critical infrastructure, this
window may exceed the safe operational horizon.

The tiered escalation model (Section~\ref{sec:bailout-protocol}) partially
mitigates this by progressively narrowing the action space at each tier, reducing
the decision burden on the human anchor.  However, the latency bound remains
non-deterministic and varies with operator state, concurrent workload, and
channel quality.  The protocol cannot eliminate latency inherent to human
cognition; it can only structure its impact.  Formal characterization of the
latency envelope across each tier, and rigorous minimum timing margins for
specific deployment contexts, remains an open problem.

% ------------------------------------------------------------------------------
\subsubsection{Adversarial Conditions and Operator Compromise}

The security assumptions underpinning the human anchor degrade under adversarial
conditions.  In a compromised network environment, bail-out signals may be
intercepted, delayed, falsified, or suppressed by an active adversary.  More
critically, if the human operator's authentication credentials or terminal are
compromised, an adversary may issue fraudulent bail-out signals that forcibly
disconnect legitimate agents, effectively weaponizing the protocol as a
denial-of-service mechanism.

The protocol currently assumes a trusted operator identity verification layer
external to the BX3 Framework.  This assumption must be validated against threat
models in which the operator's terminal \emph{is} the attack surface.  The
credential binding described in Section~\ref{sec:guarantee-accountability}
secures the resolution transition, but it does not protect the escalation signal
itself during transit.  End-to-end authenticity and integrity of the escalation
payload are necessary conditions for the protocol's security guarantees to hold
under active attack.

% ------------------------------------------------------------------------------
\subsubsection{Scalability of Human Oversight}

The current architecture designates a single human anchor per agent node.  As
agentic systems scale to hundreds or thousands of concurrent operational units,
a single operator cannot maintain meaningful oversight fidelity across the full
agent population.  Empirical studies of human monitoring tasks establish a
practical ceiling of approximately 5--7 concurrent information streams before
operator performance degrades significantly~\cite{wickens2008multiple}.

The tiered filtering of the escalation model provides a meaningful narrowing
mechanism, but the initial detection and classification burden at the monitoring
interface remains a single point of failure.  Distributed human anchor pools
(Section~\ref{sec:future_work}) are the architecturally appropriate response, but
introduce coordination overhead, consensus latency, and the possibility of
conflicting verdicts that must be resolved.  Until such pools are implemented,
the protocol's scalability guarantee is bounded by the monitoring capacity of
the designated anchor.

% ------------------------------------------------------------------------------
\subsubsection{Exploitation Risk and Protocol Manipulation}

Because the bail-out protocol defines a structured, machine-readable disengagement
mechanism, adversarial agents or malicious actors who understand its structure may
engineer false-positive escalations to consume the operator's attention and
response capacity.  A sophisticated agent could induce repeated escalations by
systematically producing ambiguous or borderline veto conditions, effectively
disabling competing agents while degrading the operator's trust in genuine
escalation signals.

This exploitation risk is compounded by the current tier-transition logic: the
deterministic mapping from confidence scores and veto states to tier transitions
is learnable by an adversary with observational access to the signal stream.  Once
the mapping is characterized, an adversarial agent can produce signals calibrated
to occupy specific tiers without crossing thresholds that would trigger immediate
human notification.  Mitigation requires stochastic elements in tier-transition
logic and anomaly detection at the operator interface to identify systematic
gaming patterns.

\subsection{Future Work}
\label{sec:future_work}

The foregoing limitations define a research agenda that proceeds in three
interdependent phases, each designed to be independently valuable while
incrementally advancing the protocol toward production-grade robustness.

% ------------------------------------------------------------------------------
\subsubsection{Formal Verification of Bail-Out Invariants}

A critical next step is the application of formal methods to prove that the
bail-out protocol maintains its safety, liveness, and accountability guarantees
under all reachable states.  We intend to model the protocol as a hybrid automata
system and apply model checking (e.g., UPPAAL or SpaceEx) to verify: (i) that a
bail-out signal is always generated within a bounded time of any veto condition
persisting beyond its tolerance window; (ii) that the tier-transition function is
acyclic and terminates; and (iii) that no reachable state permits autonomous
execution and a veto condition to be simultaneously active without escalation.

Formal verification will also characterize the protocol's completeness: the
conditions under which the protocol may fail to generate a necessary bail-out
signal, and the minimum assumptions required for the guarantees to hold.  This
analysis will produce a machine-verified safety case supplementing the proof
sketches in Section~\ref{sec:guarantees}.

% ------------------------------------------------------------------------------
\subsubsection{Adaptive Timeout Mechanisms}

Current timeout thresholds are statically configured per tier.  We propose an
adaptive timeout framework that adjusts escalation latency based on: (a)
historical operator response times weighted by time-of-day, shift pattern, and
task domain; (b) real-time assessment of network channel quality and signal
latency; and (c) the valence and magnitude of the veto condition, where higher
severity triggers shorter adaptive windows.

Machine learning models trained on historical escalation logs can predict
operator availability and dynamically compress timeouts when the probability of a
rapid response is high.  This reduces unnecessary autonomous operation during
ambiguous conditions while preserving adequate response windows when the operator
is unavailable or the veto severity is elevated.  The adaptive framework must be
formally verified to ensure that timeout compression never reduces the safety
margin below the minimum required by the deployment context.

% ------------------------------------------------------------------------------
\subsubsection{Distributed Human Anchor Pools}

To address the scalability limitation, we propose replacing the single-anchor
model with a distributed pool of human anchors, each capable of receiving,
evaluating, and responding to bail-out signals for any agent in the operational
cohort.  A consensus mechanism (e.g., threshold voting or Byzantine-fault-tolerant
agreement) determines the authoritative verdict when multiple anchors respond
concurrently.

This architecture provides fault tolerance against individual operator
unavailability, increases aggregate oversight bandwidth, and reduces the attack
surface for operator-compromise attacks.  Key implementation challenges include:
designing a low-latency consensus protocol compatible with the bail-out timeline;
managing equitable workload distribution across the pool; establishing clear
accountability when the pool reaches a split or inconclusive verdict; and
ensuring that the consensus outcome satisfies the same credential-binding and
Forensic Ledger requirements as the single-anchor resolution.  The distributed model
must preserve the ACCOUNTABILITY guarantee: the identity of every participating
anchor and the basis for the consensus verdict must be permanently recorded.

\subsection{Research Roadmap}

The three phases outlined above define a sequential research program:

\begin{enumerate}
\item \textbf{Formalization.} Apply formal verification to prove safety,
  liveness, and accountability invariants under the current single-anchor model.
  Produce a machine-verified safety case suitable for regulatory submission.
\item \textbf{Adaptation.} Deploy adaptive timeout frameworks validated against
  adversarial conditions in simulated environments.  Characterize the tradeoff
  between latency compression and safety margin reduction.
\item \textbf{Distribution.} Implement and benchmark distributed human anchor
  pools with consensus-based verdicts across large agent populations.  Verify that
  the ACCOUNTABILITY guarantee is preserved under Byzantine fault tolerance.
\end{enumerate}

Each phase produces independently deployable improvements while laying the
groundwork for the next.  The foundation is the mandatory architectural
bail-out protocol as specified: complete, formally characterized, and ready for
integration into any autonomous multi-agent system that requires structurally
defensible human accountability.
